\chapter{Tầng dữ liệu}

\section{Bốn kho lưu trữ và lý do tồn tại}

\begin{longtable}{L{2.6cm} L{4.0cm} L{7.0cm}}
\toprule
\textbf{Kho} & \textbf{Nội dung} & \textbf{Lý do chọn}\\
\midrule
\endhead
PostgreSQL & Tài khoản, khoá học, tiến độ, XP, từ vựng, trạng thái khái niệm — khoảng 70 bảng. &
Dữ liệu có quan hệ chặt, cần giao dịch (\en{transaction}) và ràng buộc toàn vẹn. Đây là \textbf{nguồn sự thật}.\\
Redis & Bộ nhớ đệm phản hồi, phiên, danh sách đen thẻ, hạn ngạch, môi giới Celery, đệm đồ thị con L0/L1. &
Đọc/ghi dưới mili-giây, cấu hình \code{allkeys-lru} nên có thể mất dữ liệu — vì vậy chỉ chứa thứ tái tạo được.\\
MongoDB & Lịch sử hội thoại, tạo tác AI, vùng đệm quan sát người học. &
Tài liệu không đồng nhất, ghi nhiều đọc ít, không cần \en{join}.\\
KuzuDB & Đồ thị tri thức: khái niệm, quan hệ tiên quyết/liên quan/dễ nhầm. &
Cơ sở dữ liệu đồ thị \textbf{nhúng} (không cần máy chủ riêng), tối ưu cho duyệt theo chiều rộng trên chuỗi tiên quyết.\\
\bottomrule
\caption{Bốn kho lưu trữ}
\end{longtable}

\begin{luuy}[Chỉ hai kho chứa dữ liệu người dùng]
Chính sách sao lưu chỉ bảo vệ PostgreSQL và MongoDB. Redis là bộ nhớ đệm và cố ý
\textbf{không} sao lưu; thư mục \code{ai\_models} và \code{content\_etl\_data} tải lại được.
\end{luuy}

\section{Nhóm bảng PostgreSQL theo miền}

\begin{longtable}{L{3.4cm} L{10.2cm}}
\toprule
\textbf{Miền} & \textbf{Bảng chính}\\
\midrule
\endhead
Người dùng & \code{users}, \code{user\_devices}, \code{refresh\_tokens}, \code{partner\_api\_keys}, \code{roles}, \code{permissions}, \code{role\_permissions}, \code{audit\_logs}, \code{user\_entitlements}.\\
Nội dung khoá học & \code{courses}, \code{units}, \code{lessons}, \code{media\_resources}, \code{lesson\_vocabulary\_items}, \code{course\_categories}, \code{content\_provenance}.\\
Ngân hàng tri thức & \code{vocabulary\_items}, \code{grammar\_items}, \code{question\_bank}, \code{test\_exams}, \code{game\_words}.\\
Tiến độ học & \code{user\_course\_progress}, \code{lesson\_completions}, \code{user\_progress}, \code{lesson\_attempts}, \code{question\_attempts}, \code{daily\_activities}, \code{streaks}.\\
Từ vựng cá nhân & \code{user\_vocabulary}, \code{user\_vocab\_knowledge}, \code{vocabulary\_reviews}, \code{vocabulary\_decks}, \code{vocabulary\_deck\_items}, \code{daily\_review\_sessions}, \code{user\_grammar\_items}.\\
Mô hình người học & \code{learner\_concept\_states}, \code{learner\_state\_profiles}, \code{learner\_observation\_events}, \code{learner\_errors}, \code{mistake\_notebook\_entries}.\\
Năng lực CEFR & \code{user\_proficiency\_profiles}, \code{user\_skill\_scores}, \code{user\_level\_history}, \code{exercise\_attempts}, \code{skill\_daily\_stats}, \code{level\_assessment\_tests}.\\
Động lực học & \code{xp\_transactions}, \code{achievements}, \code{user\_achievements}, \code{user\_wallets}, \code{wallet\_transactions}, \code{leaderboard\_entries}, \code{shop\_items}, \code{user\_inventory}, \code{challenge\_reward\_claims}, \code{user\_reward\_grants}, \code{game\_sessions}.\\
Xã hội & \code{user\_following}, \code{activity\_feeds}.\\
Thông báo & \code{notifications}, \code{notification\_campaign\_jobs}, \code{user\_reminder\_preferences}, \code{reminder\_deliveries}.\\
Tác tử AI & \code{content\_agent\_jobs}, \code{content\_agent\_uploads}, \code{ranking\_agent\_jobs}.\\
Vận hành & \code{api\_cache\_entries}, \code{product\_events}.\\
\bottomrule
\caption{Bảng PostgreSQL nhóm theo miền nghiệp vụ}
\end{longtable}

\section{Ba bảng cốt lõi của mô hình người học}

\subsection{\code{learner\_concept\_states}}

Mỗi dòng là trạng thái trí nhớ của một người học đối với một khái niệm:

\begin{longtable}{L{4.4cm} L{2.4cm} L{6.8cm}}
\toprule
\textbf{Cột} & \textbf{Kiểu} & \textbf{Ý nghĩa}\\
\midrule
\endhead
\code{user\_id}, \code{concept\_id} & UUID, chuỗi & Khoá tổ hợp xác định trạng thái.\\
\code{mastery\_probability} & \code{float} & Mức nắm vững, \([0{,}01; 0{,}99]\), mặc định 0,5.\\
\code{stability\_days} & \code{float} & Độ bền trí nhớ tính bằng ngày, \([0{,}25; 3650]\).\\
\code{difficulty} & \code{float} & Độ khó cảm nhận của khái niệm với riêng người học này, \([0;1]\).\\
\code{attempt\_count}, \code{correct\_count}, \code{error\_count} & \code{int} & Thống kê tích luỹ.\\
\code{last\_interacted\_at} & thời điểm & Mốc tính thời gian trôi qua để suy giảm trí nhớ.\\
\code{next\_review\_at} & thời điểm & Hạn ôn kế tiếp do FSRS tính.\\
\code{state\_version} & \code{int} & Tăng sau mỗi lần cập nhật, dùng cho kiểm soát tương tranh.\\
\code{algorithm\_version} & chuỗi & Hiện là \code{delta-fsrs-v3}; cho phép nhận biết dòng nào tính bằng phiên bản cũ.\\
\bottomrule
\caption{Cấu trúc \code{learner\_concept\_states}}
\end{longtable}

\subsection{\code{learner\_state\_profiles}}

Chỉ hai trường có ý nghĩa: \code{user\_id} và \code{state\_epoch}. \code{state\_epoch} là
một số nguyên tăng mỗi khi trạng thái người học thay đổi. Nó là chìa khoá làm cho bộ nhớ
đệm cá nhân hoá an toàn: khoản mục đệm nào sinh ra ở \en{epoch} cũ sẽ bị từ chối, mà
không cần lưu bất kỳ điểm nắm vững nào vào bộ đệm.

\subsection{\code{learner\_observation\_events}}

Bảng \en{outbox}. Mỗi dòng là một quan sát ("người học trả lời khái niệm X đúng/sai với
độ tin cậy $c$ vào lúc $t$"), có \code{event\_id} \textbf{duy nhất} để chống trùng.
Vòng đời được ghi thẳng vào cột \code{status}: \code{pending} $\rightarrow$
\code{claimed} $\rightarrow$ \code{applied}. Các cột \code{available\_at},
\code{claimed\_at}, \code{attempt\_count}, \code{last\_error\_code} phục vụ cơ chế thuê
(\en{lease}) và thử lại.

\begin{luuy}[Vì sao cần bảng này thay vì gọi trực tiếp]
Hội thoại AI sinh quan sát ở ai-service, còn trạng thái nằm ở backend. Nếu gọi đồng bộ,
một lần mạng chập sẽ mất bằng chứng học tập. Với \en{outbox}: ai-service ghi vào vùng đệm
MongoDB trước khi trả lời người dùng; một tiến trình thuê chuyển tiếp sang backend;
backend ghi dòng \en{outbox} \textbf{trong cùng giao dịch} rồi mới xác nhận. Nhiều tiến
trình áp dụng song song bằng \codesp{FOR UPDATE SKIP LOCKED}, nên trạng thái, \en{epoch} và
cờ đã-áp-dụng luôn thay đổi trong một giao dịch duy nhất.
\end{luuy}

\section{Dữ liệu năng lực CEFR}

\begin{itemize}[leftmargin=1.4em]
  \item \code{user\_skill\_scores} — một dòng cho mỗi (người học, kỹ năng); giữ điểm
  hiện tại, độ tin cậy, và hai cột lịch sử \code{score\_7d\_ago}/\code{score\_30d\_ago}.
  \item \code{exercise\_attempts} — \textbf{nhật ký chi tiết chỉ ghi, không đọc lại}.
  Được cắt tỉa sau 90 ngày. Không được dùng để dựng biểu đồ lịch sử.
  \item \code{skill\_daily\_stats} — thứ thực sự tồn tại lâu dài: một dòng cho mỗi
  (người học, ngày, kỹ năng), ghi bằng \en{UPSERT}. Nhờ vậy một ngày học tốn tối đa sáu
  dòng dù người học luyện tập bao nhiêu. Mọi biểu đồ lịch sử đọc từ đây.
  \item \code{user\_level\_history} — lịch sử thăng/giáng bậc CEFR.
  \item \code{level\_assessment\_tests} — bài kiểm tra xếp bậc.
\end{itemize}

\section{Ràng buộc dữ liệu không hiển nhiên}

Những ràng buộc sau đây không đọc được từ tên cột, nên đã gây lỗi ít nhất một lần:

\begin{longtable}{L{5.2cm} L{8.4cm}}
\toprule
\textbf{Trường} & \textbf{Ràng buộc}\\
\midrule
\endhead
\code{VocabularyItem.part\_of\_speech} & Kiểu liệt kê \code{PartOfSpeech}, giá trị \textbf{chữ thường}: \code{noun}, \code{verb}, \ldots\\
\code{VocabularyItem.difficulty\_level} & Kiểu liệt kê \code{DifficultyLevel}, giá trị \textbf{chữ hoa}: \code{A1}, \code{A2}, \ldots\\
\code{GrammarItem.content} & Kiểu \code{Text}, không phải JSON; riêng \code{examples} mới là danh sách JSON.\\
\code{GameWord.cefr\_level} & Chuỗi thuần \code{A1}--\code{C2}, không phải kiểu liệt kê.\\
\code{TestExam.question\_ids} & Danh sách JSON, được điền lúc chạy chứ không phải khi gieo dữ liệu.\\
\code{Lesson.exercise\_count} & Đồng bộ tự động qua \en{hook}; gán tay sẽ bị ghi đè.\\
\bottomrule
\caption{Ràng buộc dữ liệu cần nhớ}
\end{longtable}

\section{Bộ nhớ đệm Redis}

\begin{longtable}{L{4.6cm} L{4.4cm} L{4.2cm}}
\toprule
\textbf{Mục đích} & \textbf{Mẫu khoá} & \textbf{Thời gian sống}\\
\midrule
\endhead
Danh sách đen thẻ đăng xuất & \code{blacklist:\{token\}} & Đến khi thẻ hết hạn\\
Giới hạn tần suất & \code{ratelimit:\{ip\}:\{window\}} & Cửa sổ trượt\\
Đệm phản hồi API & \code{cache:\{endpoint\}:\{params\}} & Cấu hình theo điểm cuối\\
Hạn ngạch người dùng & \code{quota:\{user\_id\}:\{period\}} & Theo chu kỳ\\
Đệm TRACE-CAG lớp L0 & \code{cache\_L0:\{hash\}} & 3.600 giây\\
Đệm TRACE-CAG lớp L1 & \code{L1:\{bucket\}} & 7.200 giây\\
Hồ sơ / lịch sử hội thoại & \code{profile:\{user\_id\}}, \code{history:\{session\_id\}} & Theo phiên\\
\bottomrule
\caption{Không gian khoá Redis}
\end{longtable}

\section{MongoDB}

\begin{longtable}{L{5.4cm} L{8.2cm}}
\toprule
\textbf{Bộ sưu tập} & \textbf{Nội dung}\\
\midrule
\endhead
\code{chat\_sessions}, \code{chat\_messages} & Hội thoại thông thường theo chủ đề.\\
\code{lexi\_sessions}, \code{lexi\_messages} & Hội thoại với trợ lý Lexi chạy TRACE-CAG.\\
\code{learner\_observation\_spool} & Vùng đệm quan sát trước khi chuyển sang backend.\\
\bottomrule
\caption{Bộ sưu tập MongoDB}
\end{longtable}

Chỉ mục quan trọng: \code{session\_id} là duy nhất trên hai bảng phiên; các bảng tin nhắn
có chỉ mục tổ hợp \code{(session\_id, timestamp)}; bảng phiên có chỉ mục
\code{(user\_id, last\_activity)}.

\section{KuzuDB — đồ thị tri thức}

Lược đồ tối giản: đỉnh là \code{Concept} (mang bậc CEFR, mô tả, từ khoá), cạnh gồm
\code{prerequisite\_of} (tiên quyết), \code{related\_to} (liên quan) và
\code{confused\_with} (dễ nhầm lẫn). Đồ thị sản xuất hiện có khoảng 15.129 khái niệm và
145.387 cạnh.

\begin{luuy}[Đồ thị không chứa dữ liệu người dùng]
KuzuDB chỉ lưu \textbf{cấu trúc tri thức dùng chung}, hoàn toàn không phụ thuộc người
dùng. Trạng thái thưa của từng người học nằm ở PostgreSQL. Tách như vậy để một đồ thị
duy nhất phục vụ được mọi người học, và để bộ nhớ đệm đồ thị con có thể chia sẻ giữa các
người dùng mà không rò rỉ thông tin cá nhân.
\end{luuy}
